\chapter{Tubal Ligation}
\section*{Overview}
Tubal ligation permanently blocks, cuts, or removes the fallopian tubes to prevent pregnancy.
\section*{Common Instruments and Equipment}
The exact setup varies by technique and setting. Common items may include:
\begin{itemize}[nosep,leftmargin=*]
  \item laparoscopic camera and ports or open surgical instruments
  \item clips, rings, or electrosurgical device depending on technique
  \item specimen containers if tubes are removed
\end{itemize}

\section*{Why It Is Done}
It is an elective permanent contraceptive procedure for people who are certain they do not want future pregnancies.
\section*{How It Is Performed}
The tubes may be clipped, sealed, divided, or removed laparoscopically or during a cesarean delivery. Anesthesia and approach depend on timing and circumstances.
\section*{Contraindications}
No single contraindication applies to every patient; the treating team assesses whether the procedure is suitable for the individual and the reason for treatment.

\section*{Precautions}
The care team reviews relevant health conditions, medicines, allergies, and procedure-specific risks, and explains any changes needed for the individual.

\section*{Risks and Recovery}
Bleeding, infection, injury to nearby organs, anesthesia complications, failure with ectopic pregnancy risk, and regret are possible. It does not protect against sexually transmitted infections.
\section*{References}
\begin{enumerate}\item MedlinePlus. Female Sterilization. U.S. National Library of Medicine; 2025.\item Current Medical Diagnosis and Treatment. McGraw-Hill; 2025.\end{enumerate}
\clearpage
